%! Simplifying Rational Expressions — Unit 5, Lesson 5.1 — Warm-Up
\documentclass[10pt]{article}
\usepackage{algebra2-article}
\usepackage{algebra2-boxes}
\newcommand{\TallMath}[1]{$\displaystyle #1\rule[-1.4em]{0pt}{3.2em}$}

\begin{document}

\pageheader{Unit 5, Lesson 5.1}{Warm-Up}
\namedateperiod

\vspace{0.10in}

\begin{notesbox}{Getting Ready: Skills We'll Use Today}
\small Three skills, one per idea in today's lesson. Work quickly.

\vspace{4pt}
\textbf{1. Factor completely} (the Unit~4 toolkit --- GCF first, then count the terms).
\begin{enumerate}[label=(\alph*), leftmargin=2.1em, itemsep=3pt]
  \item $x^2-9=\blank{3.4cm}$
  \item $x^2-x-6=\blank{3.4cm}$
  \item $3x^2-12=\blank{3.8cm}$ \quad (GCF \emph{first})
\end{enumerate}

\vspace{4pt}
\textbf{2. What are you allowed to cancel?} Reduce each numerical fraction, then settle an argument.
\begin{enumerate}[label=(\alph*), leftmargin=2.1em, itemsep=3pt]
  \item Write top and bottom as products of $6$ and something: $18=6\cdot\blank{0.8cm}$ \; and \;
        $24=6\cdot\blank{0.8cm}$, \; so $\dfrac{18}{24}=\blank{1.4cm}$
  \item Kai looks at $\dfrac{3+6}{3}$, crosses out both $3$s, and writes $6$. Compute the
        \emph{true} value: $\dfrac{3+6}{3}=\blank{1.2cm}$. \; So Kai is \blank{1.6cm}.
  \item Complete the rule: you may divide out a common \blank{2.2cm}, but never a
        \blank{2.2cm}.
\end{enumerate}

\vspace{4pt}
\textbf{3. Where is the expression undefined?} (From Lesson~5.0.) For
$\dfrac{x+1}{x^2-4}$:
\begin{enumerate}[label=(\alph*), leftmargin=2.1em, itemsep=3pt]
  \item To find the excluded values, set the \blank{3.2cm} equal to \blank{1.0cm} and solve:
        $x=\blank{2.4cm}$
  \item Domain, as a union: \blank{6.6cm}
  \item The numerator is zero at $x=-1$. Is $x=-1$ an excluded value? \blank{1.2cm} \;
        Why? \blank{4.6cm}
\end{enumerate}

\end{notesbox}

\end{document}
